% Nouveau EGA V (cote 25): archive cote 25, site page 34 (PDF page 35).
% typescript. Draft compared: none.
% AI-transcribed by Claude Opus 5.5 (claude-opus-5-5, adaptive thinking, effort high), Message Batches API, batch ega25b (msgbatch_01A7n5YuXubUPNK1zsepi8RG), 2026-11-13; not yet checked by eye.
\scanpage{25/34}
% NOTE: The inserted continuation after une composante is typed at the top of the sheet and linked there by a handwritten placement stroke.
\var{\margin{30}}{\placedmargin{30}}{marginal note placed}{ega25_25_034-v0001}
Signalons le résultat dimensionnel suivant:

\uline{Proposition} 8.4. Soit T une partie fermée de X, et supposons
dodim(T,X) $\hand{\geq}$ k. Alors pour tout \hand{$\xi\in$}$\overset{\hand{\vee}}{P}$, on a dodim ($T_{\hand{\xi}}$,$Y_{\hand{\xi}}$) $\hand{\geq}$ k-1. Soit
Z l'ensemble des \hand{$\xi\in$} $\overset{\hand{\vee}}{P}$ tels que codim ($T_{\hand{\xi}}$,$Y_{\hand{\xi}}$ ) = k-1, (i.e.
codim ($T_{\hand{\xi}}$,$Y_{\hand{\xi}}$) $<$ k) alors Z est une partie constructible rare de $\overset{\hand{\vee}}{P}$ ,
\var{\inserted[typed]{i.e.}}{i.e.}{insertion}{ega25_25_034-v0002}\var{\xout{et pour que}}{}{typewriter strike-out}{ega25_25_034-v0003} $\hand{\overline{\textcolor{black}{Z}}}$ \var{\inserted[typed]{est}}{est}{insertion}{ega25_25_034-v0004}\var{\xout{soit}}{}{typewriter strike-out}{ega25_25_034-v0005} de codimension $\hand{\geq}$ 1 dans $\overset{\hand{\vee}}{P}$ . Pour qu'il soit de
\var{\deleted{\illegible{5}}}{}{deletion}{ega25_25_034-v0006}
codimension $\hand{\geq}$ 2, il f et s que pour toute composante irréductible
$T_i$ de X de codimension \var{\inserted[typed]{égale à k, \var{\xout{\unsure{telle que il|tel que il}}}{}{typewriter strike-out}{ega25_25_034-v0008}\var{\deleted{existe}}{}{deletion}{ega25_25_034-v0009} et telle}}{égale à k, \var{\xout{\unsure{telle que il|tel que il}}}{}{typewriter strike-out}{ega25_25_034-v0008}\var{\deleted{existe}}{}{deletion}{ega25_25_034-v0009} et telle}{insertion}{ega25_25_034-v0007} \var{\xout{égale à k, \unsure{on a|ont} dim f($T_i$) \illegible{1} 0, et \illegible{1} compo-}}{}{typewriter strike-out}{ega25_25_034-v0010}
que dim f($T_i$) = 0, il e\var{\xout{xi}}{}{typewriter strike-out}{ega25_25_034-v0011}xiste une composante \var{\inserted[typed]{\var{\xout{sante}}{}{typewriter strike-out}{ega25_25_034-v0013}irr $X_j$ de X telle que codim($T_i$,$X_j$)=k, et \var{\deleted{dim f($X_j$)}}{}{deletion}{ega25_25_034-v0014} = 0, \var{\inserted{\textcolor{black}{(i.e.} si f est quasi-fini \uline{et $k>0$}, \var{\deleted{que}}{}{deletion}{ega25_25_034-v0016} T n'ait pas de pt isolé x tel que $\dim_x X = k$ ).}}{\textcolor{black}{(i.e.} si f est quasi-fini \uline{et $k>0$}, \var{\deleted{que}}{}{deletion}{ega25_25_034-v0016} T n'ait pas de pt isolé x tel que $\dim_x X = k$ ).}{insertion}{ega25_25_034-v0015}}}{\var{\xout{sante}}{}{typewriter strike-out}{ega25_25_034-v0013}irr $X_j$ de X telle que codim($T_i$,$X_j$)=k, et \var{\deleted{dim f($X_j$)}}{}{deletion}{ega25_25_034-v0014} = 0, \var{\inserted{\textcolor{black}{(i.e.} si f est quasi-fini \uline{et $k>0$}, \var{\deleted{que}}{}{deletion}{ega25_25_034-v0016} T n'ait pas de pt isolé x tel que $\dim_x X = k$ ).}}{\textcolor{black}{(i.e.} si f est quasi-fini \uline{et $k>0$}, \var{\deleted{que}}{}{deletion}{ega25_25_034-v0016} T n'ait pas de pt isolé x tel que $\dim_x X = k$ ).}{insertion}{ega25_25_034-v0015}}{insertion}{ega25_25_034-v0012}
La première assertion résulte aussitôt du lemme suivant \var{\inserted[typed]{(8.4.1.a)}}{(8.4.1.a)}{insertion}{ega25_25_034-v0017}, qui est

un remords au par.5 :

\uline{Lemme} 8.4.1. Soient X préschéma loc noeth, $\underline{L}$ un Module inversible sur
X , \hand{$\phi$} une section de $\underline{L}$, Y = V(\hand{$\phi$}), T une partie fermée de X. \var{\inserted[typed]{Supposons dodim(Y,X)$\hand{\geq}$k.}}{Supposons dodim(Y,X)$\hand{\geq}$k.}{insertion}{ega25_25_034-v0018} Alors :

a) Codim(T$\hand{\cap}$Y,Y) $\hand{\geq}$ k-1.

b) Pour qu'on ait codim(T Y,Y) = k-1 i.e. codim(T$\hand{\cap}$Y,Y) $<$ k, il
f et s qu'il existe une cmmposante irréductible $T_i$ de T, contenue dans
Y, \var{\xout{et}}{}{typewriter strike-out}{ega25_25_034-v0019} telle que codim($T_i$,X) = k, et telle que pour toute composante
irréductible $X_j$ de X \var{\xout{\unsure{px|pa}}}{}{typewriter strike-out}{ega25_25_034-v0020} contenant $T_i$ et telle que dim $\underline{O}_{X_j,T_i}$ =
dim $\underline{O}_{X,T_i}$ ; (=k ), on ait $X_j \hand{\not\subset}$ Y.

La vérification de ce lemme est immédiate grâce aux généralités
de $O_{IV}$ sur la dimension. Sous les conditions de 8.4., on voit \var{\inserted[typed]{par 8.4.1. b)}}{par 8.4.1. b)}{insertion}{ega25_25_034-v0021} quels \var{\inserted[typed]{sont}}{sont}{insertion}{ega25_25_034-v0022}
les hyperplans exceptionnels $H_{\hand{\xi}}$ \var{\xout{sont xceux xpour xlesquels xil xexiste xun}}{}{typewriter strike-out}{ega25_25_034-v0023}
\var{\xout{exemple xde xcomposantes xirréd xx ($T_i$ xxxx $X_j$) xx avec xxx codim x($T_i$,$X_j$) xx = k}}{}{typewriter strike-out}{ega25_25_034-v0024}, Si
on excepte l'ensemble $Z_o$ des \hand{$\xi\in$} $\overset{\hand{\vee}}{P}$ tels qu'il existe une composante
irr R de T ou de X telle \var{\inserted[typed]{dim f(R) $>$ 0 et}}{dim f(R) $>$ 0 et}{insertion}{ega25_25_034-v0025} que $R_{\hand{\xi}}$ soit de ``dimension trop grande'',
\var{\xout{etxxqua}}{}{typewriter strike-out}{ega25_25_034-v0026} (ensemble qui est de codimension $\hand{\geq}$ 2 et par suite ne compte pas
les $H_{\hand{\xi}}$ \var{\xout{sont}}{}{typewriter strike-out}{ega25_25_034-v0027} exceptionnels sont ceux pour lesquels il existe
un $T_i$ avec codim($T_i$,X) = k , et \var{\xout{telxquexpourxtout}}{}{typewriter strike-out}{ega25_25_034-v0028} dim f(T ) = 0,

\mspage{\hand{34}}
